\documentclass[12pt,a4paper,twoside,openright,titlepage,final]{article}
\usepackage{fontspec}
\usepackage{amsmath}
\usepackage{amsfonts}
\usepackage{amssymb}
\usepackage{makeidx}
\usepackage{graphicx}
\usepackage[hidelinks,unicode=true]{hyperref}
\usepackage[spanish,es-nodecimaldot,es-lcroman,es-tabla,es-noshorthands]{babel}
\usepackage[left=3cm,right=2cm, bottom=4cm]{geometry}
\usepackage{natbib}
\usepackage{microtype}
\usepackage{ifdraft}
\usepackage{verbatim}
\usepackage[obeyDraft]{todonotes}
\ifdraft{
	\usepackage{draftwatermark}
	\SetWatermarkText{BORRADOR}
	\SetWatermarkScale{0.7}
	\SetWatermarkColor{red}
}{}
\usepackage{booktabs}
\usepackage{longtable}
\usepackage{calc}
\usepackage{array}
\usepackage{caption}
\usepackage{subfigure}
\usepackage{footnote}
\usepackage{url}
\usepackage{tikz}
\usepackage{pdflscape}
\usepackage{minted}

%\setsansfont[Ligatures=TeX]{texgyreadventor}
%\setmainfont[Ligatures=TeX]{texgyrepagella}

%*******************************************************
%                 NO MODIFICAR
\newcommand*{\FSfont}[1]{%
  \fontencoding{T1}\fontfamily{#1}\selectfont}

\newlength{\tpheight}\setlength{\tpheight}{0.9\textheight}
\newlength{\txtheight}\setlength{\txtheight}{0.9\tpheight}
\newlength{\tpwidth}\setlength{\tpwidth}{0.9\textwidth}
\newlength{\txtwidth}\setlength{\txtwidth}{0.9\tpwidth}
\newlength{\drop}
%*******************************************************

% Crea una portada con los siguientes parámetros
%
% #1 : Título 
% #2 : Subtítulo
% #3 : Subsubtítulo
% #4 : Autor(es)
% #5 : Lugar
%

\newcommand*{\portada}[5]{
\begin{titlepage}
\begingroup
\vspace*{1cm}
\drop = 0.2\txtheight
\centering
\vfill
{\Huge \scshape #1}\\[\baselineskip]
{\Large \textbf{#2}}\\[\baselineskip]
{\Large \scshape #3}\\[\baselineskip]
\vspace*{0.3cm}
{\large \textit{#4}}\\[0.5\drop]
\includegraphics[scale=0.35]{./imagenes/logoURJC.jpg}
\vspace*{1.5cm}

{\large \scshape #5, \today} \par
\begin{center}
\end{center}
\vfill\null
\endgroup
\end{titlepage}
}
 %*****************************************************
 


\author{José Ignacio Escribano}

\setlength{\parindent}{0pt}

\begin{document}

\pagenumbering{alph}
\setcounter{page}{1}

\portada{Práctica II}{Simulación y Metaheurísticas}{Búsquedas locales}{José Ignacio Escribano}{Móstoles}

\listoftables
\thispagestyle{empty}
\newpage

\listoffigures
\thispagestyle{empty}
\newpage

\tableofcontents
\thispagestyle{empty}
\newpage


\pagenumbering{arabic}
\setcounter{page}{1}

\section{Introducción}

En esta práctica realizaremos búsquedas locales para el problema del p-hub.

\section{Resolución de la práctica}

A continuación, resolveremos todos las cuestiones planteadas de la práctica.

\subsection{Definir una búsqueda local}

El código para realizar una búsqueda local a partir de su vecindad y el tipo de orden (lexicográfico o aleatorio) se puede ver a continuación (función \texttt{busquedaLocal} de la clase \texttt{Utils}):

\begin{minted}[linenos=true, numberblanklines=true, breaklines=true]{java}
public static Solución busquedaLocal(String tipo_orden, List<Solución> vecindades, Solución actual) {
    Solución mejor_solución = actual;

    if (tipo_orden == "lexicográfico") {
        for (int i = 0; i < vecindades.size(); i++) {
            if (vecindades.get(i).getObjetivo() < mejor_solución.getObjetivo()) {
                mejor_solución = vecindades.get(i);
            }
        }
    } else if (tipo_orden == "aleatorio") {
        // Generamos un número aleatorio entre 0 y número de vecindades

        // Copiamos la lista con las vecindades
        List<Solución> vecindades_copia = new ArrayList<Solución>(vecindades);

        int n = vecindades_copia.size();

        Solución s = null;

        while (n > 1) {
            n = vecindades_copia.size();
            Random r = new Random();
            int indice = r.nextInt(n);
            s = vecindades_copia.get(indice);

            if (s.getObjetivo() < mejor_solución.getObjetivo()) {
                mejor_solución = s;
            }

            // Eliminamos de la vecindad
            vecindades_copia.remove(s);
        }

    }
    return mejor_solución;
}
\end{minted} 

En las líneas 4 a 9 se realiza la búsqueda por orden lexicográfico y de las líneas 10 a 32 se realiza por orden aleatorio.\\
En el caso de la búsqueda lexicográfica se comparan las soluciones de la vecindad con la mejor solución actual, se selecciona ésta como mejor solución actual y se continúa el proceso hasta que se acaban las soluciones de las vecindad.\\
Notar que seguimos el orden que viene en la lista de soluciones vecinas, puesto que éstas vienen en ese orden.\\
En el caso del orden aleatorio, se genera un número aleatorio en el intervalo discreto $[0, \text{número de elementos vecindad})$ y se selecciona el índice que corresponde con el número aleatorio, se compara con la mejor solución actual y se elimina de la vecindad.\\

El código para obtener la vecindad a partir de una solución y su instancia está en la función \texttt{generarVecindad} de la clase \texttt{PHub}. Esta función busca un cliente y un servidor, intercambia sus papeles y genera una nueva matriz de adyacencia con esa configuración.

\begin{minted}[linenos=true, numberblanklines=true, breaklines=true]{java}
static List<Solución> generarVecindad(Solución s, InstanciaPHub instancia) {
    List<Solución> vecindad = new ArrayList<>();
    boolean[] sol = s.getSolucion();
    boolean[] sol2 = new boolean[sol.length];

    int nodos = sol.length;

    // Generamos las permutaciones
    for (int i = 0; i < sol.length; i++) {
        for (int j = 0; j < sol.length; j++) {
            // Buscamos un true y un falso entre un servidor y un cliente
            if (sol[i] == true && sol[j] == false) {
                // Copiamos el array sol
                System.arraycopy(sol, 0, sol2, 0, sol.length);

                // Intercambiamos las posiciones
                sol2[i] = false;
                sol2[j] = true;

                // Creamos la matriz de adyacencia de la nueva solución
                boolean[][] ady = new boolean[nodos][nodos];
                for (int z = 0; z < nodos; z++) {
                    if (!sol2[z]) {
                        // Seleccionamos el servidor más cercano que nos
                        // encontremos
                        int serv = Utils.seleccionarServidor(z, sol2, instancia.getDistancia());

                        ady[z][serv] = true;
                        ady[serv][z] = true;
                     }

                }

                Solución s1 = new Solución(sol2, ady, instancia.getDistancia());

                if (Utils.esSoluciónVálida(s1, instancia)) {
                    vecindad.add(s1);
                }
            }
        }
    }
    return vecindad;
}
\end{minted}

En las líneas 9 a 12 se busca un servidor y un cliente y se intercambian sus posiciones. De la línea 21 a la 32, se genera la matriz de adyacencia de la misma forma que para generar una solución aleatoria. Se crea la solución y se comprueba si es válida. En caso afirmativo se añade a la lista de vecinos. Por último, se devuelve la lista de vecinos.

\subsection{Programación de los dos algoritmos}

Elegimos la construcción aleatoria y mejora con orden de exploración lexicográfico (A2) y orden de exploración aleatorio (A3). El código de estos dos algoritmos se implementa en la clase \texttt{Práctica2}.

\begin{minted}[linenos=true, numberblanklines=true, breaklines=true]{java}
public static Solución constAleatoriaYMejoraLexicográfica(Solución sol, InstanciaPHub instancia){

    Solución result = null;
    List<Solución> vecindad = PHub.generarVecindad(sol, instancia);
    result = Utils.busquedaLocal("lexicográfico", vecindad, sol);
    return result;
}

public static Solución constAleatoriaYMejoraAleatorio(Solución sol, InstanciaPHub instancia){

    Solución result = null;
    List<Solución> vecindad = PHub.generarVecindad(sol, instancia);
    result = Utils.busquedaLocal("aleatorio", vecindad, sol);
    return result;
}
\end{minted}

En ambos casos se genera la vecindad a partir de la solución que se pasa como parámetro, se genera la vecindad, y se realiza la búsqueda según el tipo (lexicográfico o aleatorio).


\subsection{Análisis de los resultados}


\section{Conclusiones}


\end{document}